% Source: OpenAI2026AreaLaw, 09-amplification.tex, lines 304--322
% (proof of prop:amplification, eq:amplification-typical-error), with
% Lemma 3.1 (lem:tail), 02-initial.tex, lines 34--69 and 202--206,
% revision openai/math@adc7f1241b42e322a6451854ab7e4b4c146bf78a.

\section{Typical projectors for the amplification}

The proof of the amplification proposition in
\cite[Section~10]{OpenAI2026AreaLaw} begins with three spectral
projections on the cuts $X$, $Y$, and $S_0$. We construct these projections
from the marginal tail bound and bound the failure probability of each.

\begin{definition}[Typical projection of a cut]
    \label{def:al_amp_typical_projection}
    \lean{TNLean.PEPS.AreaLaw.typicalProjection}
    \leanok
    \uses{def:area_law_ground_entropy}
    Let $\Omega$ be a vector on a finite induced lattice domain $\Lambda$,
    let $A\subseteq\Lambda$, and let $w\in\mathbb R$. Write
    $\rho_A=\sum_jp_j|e_j\rangle\langle e_j|$ for an eigendecomposition
    of the reduced state $\rho_A$ of $\Omega$ on $A$. The typical projection
    of $A$ at width $w$ is
    \begin{align*}
        P_A^w=\sum_{j:\,p_j>0,\ |-\log p_j-S_\Omega(A)|\le w}
        |e_j\rangle\langle e_j|.
    \end{align*}
\end{definition}

\begin{theorem}[Typical projection of one cut]
    \label{thm:al_amp_typical_projection_cut}
    \lean{TNLean.PEPS.AreaLaw.typicalProjection_failure_le_surprisalTail}
    \leanok
    \uses{def:al_amp_typical_projection, thm:area_law_marginal_tails}
    Let $\Omega$ be a unit vector, $A\subseteq\Lambda$, and $w\in\mathbb R$.
    Then $P_A^w$ is an orthogonal projection commuting with $\rho_A$. Every
    retained eigenvalue satisfies $p_j>0$ and
    $\mathrm e^{-S_\Omega(A)-w}\le p_j\le\mathrm e^{-S_\Omega(A)+w}$.
    With $K$ as in Theorem~\ref{thm:area_law_marginal_tails},
    \begin{align*}
        1-\operatorname{tr}(\rho_AP_A^w)\le\Pr\{|K-S_\Omega(A)|>w\}.
    \end{align*}
\end{theorem}
\begin{proof}\leanok
    Since $P_A^w$ is diagonal in the eigenbasis of $\rho_A$, the trace
    $\operatorname{tr}(\rho_AP_A^w)$ is the sum of the retained
    eigenvalues. Because $\operatorname{tr}\rho_A=1$ and an eigenvalue
    contributes to $1-\operatorname{tr}(\rho_AP_A^w)$ only when $p_j>0$
    and $|-\log p_j-S_\Omega(A)|>w$, the stated tail bound holds.
\end{proof}

\begin{theorem}[Failure probability of a typical projection]
    \label{thm:al_amp_typical_failure}
    \lean{TNLean.PEPS.AreaLaw.exists_typicalProjection_failure_le_stretchedExp}
    \leanok
    \uses{def:al_amp_typical_projection}
    Fix $q\ge1$, a range $R$, a strength $J\ge0$, a gap $\Delta>0$, and
    $C_0\ge0$. There exist constants $C,c>0$ with the following property.
    Let $n\ge1$, let $H$ be a range-$R$ interaction of strength $J$ on a
    finite induced lattice domain $\Lambda$, and let $\Omega$ be a unit
    ground vector with full-system gap $\Delta$. Then every cut
    $A\subseteq\Lambda$ with $|\partial_\Lambda A|\le C_0n$ satisfies
    \begin{align*}
        1-\operatorname{tr}\bigl(\rho_AP_A^{n^{3/5}}\bigr)
        \le C\exp(-cn^{1/10}).
    \end{align*}
\end{theorem}
\begin{proof}\leanok
    \uses{thm:area_law_marginal_tails, thm:area_law_tail_parameter,
        thm:al_amp_typical_projection_cut}
    Let $C_1$ be the constant of Theorem~\ref{thm:area_law_tail_parameter}
    and set $\vartheta=J/\Delta$. Then
    $\mathcal B_A\le C_1(1+C_0n)\le C_1(1+C_0)n$. By
    Theorem~\ref{thm:al_amp_typical_projection_cut} and
    (\ref{eq:area_law_marginal_tail}) at $w=n^{3/5}$, the failure
    probability is at most
    $2\mathrm e^{\mathrm e/2}\exp(-n^{3/5}/(32\sqrt{(1+\vartheta)\mathcal B_A}))$.
    Since $n^{3/5}/\sqrt{n}=n^{1/10}$, this bound is at most
    $2\mathrm e^{\mathrm e/2}\exp(-cn^{1/10})$ with
    $c=1/(32\sqrt{(1+\vartheta)C_1(1+C_0)})$.
\end{proof}

\begin{theorem}[Typical projectors of the amplification]
    \label{thm:al_amp_typical_projectors}
    \lean{TNLean.PEPS.AreaLaw.exists_typicalProjection_failure_le_inv_pow,
        TNLean.PEPS.AreaLaw.exists_typicalProjectors_three_cuts}
    \leanok
    \uses{def:al_amp_typical_projection}
    Fix $q$, $R$, $J$, $\Delta$, and $C_0$ as in
    Theorem~\ref{thm:al_amp_typical_failure}, and an integer $K\ge0$.
    There exists a threshold $N$, depending only on these parameters, such
    that for every $n\ge N$, every such $\Lambda$, $H$, and $\Omega$, and
    every three cuts $X$, $Y$, $S_0$ with at most $C_0n$ crossing edges
    each, the projections $P_X=P_X^w$, $T_Y=P_Y^w$, and $Q_{S_0}=P_{S_0}^w$
    at width $w=n^{3/5}$ are orthogonal projections commuting with their
    respective reduced states, retain only positive eigenvalues in
    $[\mathrm e^{-S_\Omega(A)-w},\mathrm e^{-S_\Omega(A)+w}]$ for each cut
    $A\in\{X,Y,S_0\}$, and each satisfy a failure bound of at most
    \begin{align*}
        \delta=n^{-2K}.
    \end{align*}
    This is the projector construction in \cite[Section~10]{OpenAI2026AreaLaw}
    preceding the Schmidt expansion of $\Omega$.
\end{theorem}
\begin{proof}\leanok
    \uses{thm:al_amp_typical_failure, thm:al_amp_typical_projection_cut}
    Since $x^{20K}\mathrm e^{-cx}\to0$ as $x\to\infty$, setting $x=n^{1/10}$
    shows that the bound $C\exp(-cn^{1/10})$ of
    Theorem~\ref{thm:al_amp_typical_failure} is at most $n^{-2K}$ for all
    sufficiently large $n$. Applying
    Theorem~\ref{thm:al_amp_typical_projection_cut} to each of the three cuts
    gives the result.
\end{proof}

% Source: OpenAI2026AreaLaw, 09-amplification.tex, lines 304--447
% (proof of prop:amplification: eq:amplification-typical-error through
% eq:amplification-ground-component), revision
% openai/math@adc7f1241b42e322a6451854ab7e4b4c146bf78a.

\section{The restoring operator of the amplification}

In \cite[Section~10]{OpenAI2026AreaLaw}, the proof of the amplification
proposition constructs a restoring operator from the typical projections and
establishes three estimates before introducing the constraint products. This
section proves those three estimates for the ground state.

Throughout, $\Omega$ is a unit vector on a finite induced lattice domain
$\Lambda$ with $q$ levels per site, $X,Y\subseteq\Lambda$ are disjoint,
$S_0=X\cup Y$, and $Z'=\Lambda\setminus S_0$. For a width $w$, write
$P_X=P_X^w$, $T_Y=P_Y^w$, and $Q=P_{S_0}^w$ for the typical projections of
Definition~\ref{def:al_amp_typical_projection}. Write
$\rho_X=\sum_xp_x|x\rangle\langle x|$ in a full orthonormal eigenbasis
$\{|x\rangle\}$, and call $x$ selected when $|x\rangle$ lies in the range
of $P_X$. Append two auxiliary systems $s',e'$, each a copy of the state space
on $X$ with a fixed unit reference vector $|0\rangle$. Define the restoring
operator $R_0$ on $s',X,Y$ and the column operator $B_0$ on $s',e',X,Y$ by
\begin{align*}
    R_0 & =\sum_{x\ \mathrm{selected}}p_x^{-1/2}
    |x\rangle\langle0|_{s'}\otimes
    Q\bigl(|x\rangle\langle0|_X\otimes T_Y\bigr),          \\
    B_0 & =\sum_{x\ \mathrm{selected}}\ \sum_{y}p_x^{-1/2}
    |x\rangle\langle0|_{s'}\otimes|y\rangle\langle0|_{e'}\otimes
    Q\bigl(|x\rangle\langle y|_X\otimes T_Y\bigr),
\end{align*}
where $y$ runs over the full eigenbasis. Let $\phi$ be $\Omega$ with $X$
relabelled as $e'$ and $|0\rangle$ tensored on $X$ and $s'$, and let
$\eta=\Omega\otimes\sum_x\sqrt{p_x}|x\rangle_{s'}|x\rangle_{e'}$.

\begin{definition}[Coordinates of the restoring operator]
    \label{def:al_amp_restoring_coordinates}
    \lean{TNLean.PEPS.AreaLaw.restoringVector,
        TNLean.PEPS.AreaLaw.restoringCutProjection}
    \leanok
    \uses{def:al_amp_typical_projection}
    Decompose the configuration space of $\Lambda$ as the product of the
    configuration spaces of $X$, $Y$, and $Z'$, and express the factor on $X$
    in the eigenbasis $\{|x\rangle\}$ of $\rho_X$. In these coordinates,
    $\Omega$ and $Q$ determine $R_0$, $B_0$, $\phi$, and $\eta$.
\end{definition}

\begin{theorem}[Mutual information of the collar]
    \label{thm:al_amp_restoring_purity}
    \lean{TNLean.PEPS.AreaLaw.mutualInformation_compl_eq}
    \leanok
    \uses{def:area_law_ground_entropy}
    For every unit vector $\Omega$ and disjoint subsets $X,Y\subseteq\Lambda$,
    \begin{align*}
        I_\Omega(X:Z')=S_\Omega(X)+S_\Omega(S_0)-S_\Omega(Y).
    \end{align*}
\end{theorem}
\begin{proof}\leanok
    Complementary regions of a pure state have equal entropy, so
    $S_\Omega(Z')=S_\Omega(S_0)$ and $S_\Omega(X\cup Z')=S_\Omega(Y)$.
\end{proof}

\begin{theorem}[Spectral inequalities of the restoring operator]
    \label{thm:al_amp_restoring_spectral}
    \lean{TNLean.PEPS.AreaLaw.restoringCutProjection_le_and_typicalProjection_le,
        TNLean.PEPS.AreaLaw.restoringCutProjection_isStarProjection_commute,
        Matrix.PosSemidef.spectralSelection_le_smul,
        Matrix.IsHermitian.spectralSelection_mul_mul_le_smul}
    \leanok
    \uses{def:al_amp_restoring_coordinates}
    The projection $Q$ commutes with $\rho_{S_0}$, and
    \begin{align*}
        Q\le\mathrm e^{S_\Omega(S_0)+w}\rho_{S_0},\qquad
        T_Y\rho_YT_Y\le\mathrm e^{-S_\Omega(Y)+w}T_Y.
    \end{align*}
\end{theorem}
\begin{proof}\leanok
    \uses{thm:al_amp_typical_projection_cut}
    Both $Q$ and $T_Y$ are spectral projections of the corresponding reduced
    states. Every eigenvalue selected by $Q$ is at least
    $\mathrm e^{-S_\Omega(S_0)-w}$, and every eigenvalue selected by $T_Y$ is
    at most $\mathrm e^{-S_\Omega(Y)+w}$. Both operator inequalities hold
    eigenvalue by eigenvalue in a common eigenbasis.
\end{proof}

\begin{theorem}[Action of the three typical projections on the ground state]
    \label{thm:al_amp_restoring_failure}
    \lean{TNLean.PEPS.AreaLaw.norm_selectedProjection_restoringVector_sub_le,
        TNLean.PEPS.AreaLaw.norm_typicalProjection_restoringVector_sub_le,
        TNLean.PEPS.AreaLaw.norm_restoringCutProjection_restoringVector_sub_le,
        Matrix.norm_projection_mulVec_sub_sq}
    \leanok
    \uses{def:al_amp_restoring_coordinates}
    If $1-\operatorname{tr}(\rho_AP_A^w)\le\delta$ for $A=X,Y,S_0$, then each
    of $P_X$, $T_Y$, and $Q$ moves $\Omega$ by at most $\sqrt\delta$ in norm.
\end{theorem}
\begin{proof}\leanok
    \uses{thm:al_amp_typical_projection_cut}
    For any orthogonal projection $P$ acting on a subsystem $A$,
    $\|P\Omega-\Omega\|^2=1-\operatorname{tr}(\rho_AP)$.
\end{proof}

\begin{theorem}[Three estimates for the restoring operator]
    \label{thm:al_amp_restoring_estimates}
    \lean{TNLean.PEPS.AreaLaw.restoringOperator_estimates}
    \leanok
    \uses{def:al_amp_restoring_coordinates}
    Suppose $1-\operatorname{tr}(\rho_AP_A^w)\le\delta$ for $A=X,Y,S_0$, and
    set $\mathcal I=I_\Omega(X:Z')$ and
    $M_0=\exp\bigl((\mathcal I+3w)/2\bigr)$. Then
    \begin{align*}
        \|R_0\|,\ \|B_0\|                    & \le M_0, \\
        \|P_\Omega R_0\phi\|                 & \le1,    \\
        \|R_0^\dagger(L')^\dagger\eta-\phi\| & \le
        3\sqrt\delta+\|(L')^\dagger\Omega-\Omega\|,
    \end{align*}
    where the third inequality holds for every physical operator $L'$.
\end{theorem}
\begin{proof}\leanok
    \uses{thm:al_amp_restoring_purity, thm:al_amp_restoring_spectral,
        thm:al_amp_restoring_failure}
    Selected eigenvalues satisfy $p_x^{-1}\le\mathrm e^{S_\Omega(X)+w}$. By
    Theorem~\ref{thm:al_amp_restoring_spectral}, the Gram
    operators of $R_0$ and $B_0$ are bounded above by
    $\mathrm e^{S_\Omega(X)+S_\Omega(S_0)-S_\Omega(Y)+3w}=M_0^2$, where the
    equality uses Theorem~\ref{thm:al_amp_restoring_purity}. To bound
    $\|P_\Omega R_0\phi\|$, set
    $M_X=\operatorname{tr}_Y[(I_X\otimes T_Y)\rho_{S_0}Q]$; since $Q$ commutes
    with $\rho_{S_0}$, we have $0\le M_X\le\rho_X$, and $\rho_X$ is diagonal in
    the basis $\{|x\rangle\}$. The adjoint error bound follows by inserting
    $P_X$, $T_Y$, and $Q$ in turn on $\Omega$ and applying
    Theorem~\ref{thm:al_amp_restoring_failure}.
\end{proof}

\begin{theorem}[The restoring operator at large scales]
    \label{thm:al_amp_restoring_large_scales}
    \lean{TNLean.PEPS.AreaLaw.exists_restoringOperator_estimates}
    \leanok
    \uses{thm:al_amp_restoring_estimates}
    Fix $q\ge1$, a range $R$, a strength $J\ge0$, a gap $\Delta>0$, $C_0\ge0$,
    and an integer $K\ge0$. There exists a threshold $N$, depending only on these
    parameters, such that for every $n\ge N$, every
    range-$R$ interaction of strength $J$ on a finite induced lattice domain
    with a unit ground vector $\Omega$ of full-system gap $\Delta$, and every
    pair of disjoint subsets $X,Y$ such that $X$, $Y$, and $S_0$ each have at
    most $C_0n$ crossing edges, the estimates of
    Theorem~\ref{thm:al_amp_restoring_estimates} hold at width $w=n^{3/5}$ with
    $3\sqrt\delta=3n^{-K}$.
\end{theorem}
\begin{proof}\leanok
    \uses{thm:al_amp_typical_projectors}
    Theorem~\ref{thm:al_amp_typical_projectors} bounds the three tail weights
    $1-\operatorname{tr}(\rho_AP_A^w)$ by $\delta=n^{-2K}$, and
    $\sqrt{n^{-2K}}=n^{-K}$.
\end{proof}

